% Ad Atticum 7.24 --- headnote
% ad-atticum-07-24, 10 February 49 BC, Formian villa

\textit{Cicero to Atticus, written from the Formian
villa on the fourth day before the Ides of February
49 BC (the manuscript dateline:
\textit{Scr. in Formiano iv Id. Febr. a. 705 (49)}).
A single short paragraph, the morning after \emph{Ad
Atticum} 7.23, walking back the rumour Philotimus had
sent. Cassius has had a letter from his friend
Lucretius at Capua relaying news from Domitius's
camp: Vibullius is hurrying with a handful of men to
Pompey, Caesar is hot on his heels, Domitius does not
have three thousand soldiers, the consuls have left
Capua.}

\textit{The arithmetic is the whole letter. ``I have
no doubt that Gnaeus is in flight; may he only get
away'' --- \emph{modo effugiat} --- is a short prayer
between two clauses. The last sentence is the
decision the previous day's news had already made:
\emph{ego a consilio fugiendi, ut tu censes, absum},
``As for me, I am keeping my distance from the plan
of flight, as you advise.'' He is still on the coast,
still at Formiae, still --- for now --- not on a
ship.}
